\chapter{Evaluación de aceleradores criptográficos poscuánticos}
\label{anexo:aceleradores}
\label{app:accelerators-evaluation}

Este anexo recoge el estudio comparativo previo a la selección del acelerador. Se conservan las valoraciones y la clasificación del documento de trabajo \textit{Accelerators Evaluation}; las tablas se han transcrito al formato de la memoria. Las valoraciones corresponden al estudio inicial y no constituyen una evaluación actualizada de la seguridad de los esquemas.

Las tablas procedentes de imágenes mantienen las cifras y abreviaturas del documento de partida. Las referencias numéricas internas de esas tablas pertenecen a sus publicaciones de origen y no a la bibliografía de esta memoria. Los guiones indican datos no disponibles en las tablas originales.


\section{CROSS}

\subsection*{Ventajas}
\begin{itemize}
\item Se compone principalmente de operaciones aritméticas que es relativamente sencillo en HW.
\item Muy paralelizable
\end{itemize}

\subsection*{Limitaciones}
\begin{itemize}
\item Tiene un gran coste en área
\item Tiene una lógica bastante compleja por lo que es necesario definir bien las FSM.
\item Mucho hashing
\end{itemize}

\begingroup
\small
\setlength{\tabcolsep}{3pt}
\setlength{\LTcapwidth}{\textwidth}
\begin{longtable}{lrrrrrr}
\caption{CROSS: recursos y frecuencia de síntesis.}\label{tab:eval-cross-recursos}\\
\toprule
Parámetros & LUT & FF & BRAM & DSP & KeSlice & MHz \\
\midrule
\endfirsthead
\toprule
Parámetros & LUT & FF & BRAM & DSP & KeSlice & MHz \\
\midrule
\endhead
\midrule
\multicolumn{7}{r}{\small Continúa en la página siguiente.}\\
\endfoot
\bottomrule
\endlastfoot
{RSDP-1-f} & {26741} & {11678} & {45.0} & {0} & {16.2} & {118} \\
{RSDP-1-b} & {27308} & {11820} & {57.5} & {0} & {19.0} & {111} \\
{RSDP-1-s} & {27779} & {11990} & {111.5} & {0} & {30.6} & {114} \\
{RSDPG-1-f} & {27324} & {11726} & {28.5} & {0} & {12.9} & {122} \\
{RSDPG-1-b} & {28572} & {11892} & {37.0} & {0} & {15.0} & {120} \\
{RSDPG-1-s} & {28777} & {12257} & {63.0} & {0} & {20.6} & {117} \\
{RSDP-3-f} & {29096} & {13761} & {97.0} & {0} & {27.8} & {108} \\
{RSDP-3-b} & {29790} & {14028} & {121.5} & {0} & {33.2} & {113} \\
{RSDP-3-s} & {29742} & {14067} & {148.0} & {0} & {38.8} & {109} \\
{RSDPG-3-f} & {31185} & {13272} & {43.5} & {0} & {17.0} & {112} \\
{RSDPG-3-b} & {31233} & {13479} & {68.0} & {0} & {22.2} & {115} \\
{RSDPG-3-s} & {32123} & {13808} & {122.5} & {0} & {34.0} & {118} \\
{RSDP-5-f} & {32085} & {15537} & {151.0} & {0} & {40.0} & {113} \\
{RSDP-5-b} & {33056} & {15635} & {202.0} & {0} & {51.1} & {111} \\
{RSDP-5-s} & {34868} & {15815} & {280.0} & {0} & {68.1} & {105} \\
{RSDPG-5-f} & {34266} & {15331} & {91.5} & {0} & {28.1} & {113} \\
{RSDPG-5-b} & {36033} & {15502} & {122.5} & {0} & {35.0} & {113} \\
{RSDPG-5-s} & {35296} & {15853} & {182.5} & {0} & {47.5} & {113} \\
\end{longtable}
\endgroup
\begingroup
\small
\setlength{\tabcolsep}{3pt}
\setlength{\LTcapwidth}{\textwidth}
\begin{longtable}{lrrrrrr}
\caption{CROSS: tiempos de generación de claves, firma y verificación (ms) y producto área-tiempo (AT).}\label{tab:eval-cross-tiempos}\\
\toprule
Parámetros & Claves & AT & Firma & AT & Verif. & AT \\
\midrule
\endfirsthead
\toprule
Parámetros & Claves & AT & Firma & AT & Verif. & AT \\
\midrule
\endhead
\midrule
\multicolumn{7}{r}{\small Continúa en la página siguiente.}\\
\endfoot
\bottomrule
\endlastfoot
{RSDP-1-f} & {0.035} & {0.568} & {0.478} & {7.76} & {0.409} & {6.64} \\
{RSDP-1-b} & {0.037} & {0.708} & {0.856} & {16.3} & {0.618} & {11.8} \\
{RSDP-1-s} & {0.036} & {1.11} & {1.67} & {51.0} & {1.12} & {34.2} \\
{RSDPG-1-f} & {0.008} & {0.109} & {0.338} & {4.35} & {0.303} & {3.91} \\
{RSDPG-1-b} & {0.009} & {0.129} & {0.676} & {10.1} & {0.519} & {7.78} \\
{RSDPG-1-s} & {0.009} & {0.182} & {1.38} & {28.4} & {1.03} & {21.1} \\
{RSDP-3-f} & {0.081} & {2.26} & {1.09} & {30.3} & {0.973} & {27.1} \\
{RSDP-3-b} & {0.078} & {2.57} & {1.69} & {56.0} & {1.24} & {41.1} \\
{RSDP-3-s} & {0.080} & {3.12} & {2.61} & {101} & {1.80} & {69.7} \\
{RSDPG-3-f} & {0.017} & {0.297} & {0.701} & {11.9} & {0.665} & {11.3} \\
{RSDPG-3-b} & {0.017} & {0.378} & {0.919} & {20.4} & {0.786} & {17.5} \\
{RSDPG-3-s} & {0.017} & {0.564} & {1.63} & {57.2} & {1.35} & {45.7} \\
{RSDP-5-f} & {0.138} & {5.51} & {1.71} & {68.4} & {1.60} & {64.1} \\
{RSDP-5-b} & {0.140} & {7.16} & {2.73} & {140} & {2.09} & {107} \\
{RSDP-5-s} & {0.148} & {10.1} & {4.62} & {314} & {3.27} & {223} \\
{RSDPG-5-f} & {0.028} & {0.798} & {1.20} & {33.6} & {1.14} & {31.9} \\
{RSDPG-5-b} & {0.028} & {0.995} & {1.57} & {54.9} & {1.32} & {46.1} \\
{RSDPG-5-s} & {0.028} & {1.35} & {2.77} & {132} & {2.15} & {102} \\
\end{longtable}
\endgroup


\section{LESS}

\subsection*{Ventajas}
\begin{itemize}
\item Hay implementaciones existentes en FPGA similares
\end{itemize}

\subsection*{Limitaciones}
\begin{itemize}
\item Se necesita mucha optimización para implementarse (alto uso de LUT).
\item Gastan mucho las operaciones vectoriales y matriciales.
\item Claves públicas muy grandes
\item Eliminación Gaussiana ocupa mucho.
\end{itemize}

\begingroup
\small
\setlength{\tabcolsep}{3pt}
\setlength{\LTcapwidth}{\textwidth}
\begin{longtable}{lrrrrrr}
\caption{LESS: recursos y prestaciones. LUT, FF y ciclos expresados en miles; BRAM de 36 Kbit.}\label{tab:eval-less}\\
\toprule
Parámetros & MHz & LUT & FF & BRAM & Ciclos & Latencia (µs) \\
\midrule
\endfirsthead
\toprule
Parámetros & MHz & LUT & FF & BRAM & Ciclos & Latencia (µs) \\
\midrule
\endhead
\midrule
\multicolumn{7}{r}{\small Continúa en la página siguiente.}\\
\endfoot
\bottomrule
\endlastfoot
{LESS-1\{b,i,s\}} & {125} & {29.4} & {20.4} & {25.5} & {16.1} & {129.3} \\
{LESS-3\{b,s\}} & {111} & {46.3} & {33.8} & {40} & {40.4} & {364.3} \\
{LESS-5\{b,s\}} & {100} & {66.6} & {47.4} & {55} & {75.7} & {756.6} \\
\end{longtable}
\endgroup


\section{SQIsign}

\subsection*{Limitaciones}
\begin{itemize}
\item La isogenía consta de un algoritmo complejo con control complejo por lo que es complicado de llevar a cabo en una FPGA.
\item Requiere mucho hardware
\end{itemize}

\begingroup
\small
\setlength{\tabcolsep}{3pt}
\setlength{\LTcapwidth}{\textwidth}
\begin{longtable}{lrrrr}
\caption{Comparación de implementaciones recogida en el estudio de SQIsign: frecuencia, latencia y recursos. Ancho en bits; latencia en ciclos.}\label{tab:eval-sqisign}\\
\toprule
Bits & Diseño & MHz & Ciclos & DSP/LUT/FF \\
\midrule
\endfirsthead
\toprule
Bits & Diseño & MHz & Ciclos & DSP/LUT/FF \\
\midrule
\endhead
\midrule
\multicolumn{5}{r}{\small Continúa en la página siguiente.}\\
\endfoot
\bottomrule
\endlastfoot
{256} & {CA0D2C1E [23]} & {625} & {140} & {16/1759/3365} \\
{256} & {[8]} & {200} & {33} & {120/2728/-} \\
{256} & {[8]} & {194} & {47} & {91/1718/-} \\
{256} & {[29]} & {--} & {--} & {289/87622/2952} \\
{256} & {[5]} & {--} & {--} & {248/9450/-} \\
{256} & {[18]} & {86.79} & {29} & {120/6688/5163} \\
{256} & {[20]} & {68} & {--} & {0/187900/-} \\
{256} & {N=3, s=6} & {625} & {111} & {18/4156/5145} \\
{256} & {N=5, s=4} & {625} & {113} & {20/4824/5541} \\
{256} & {N=7, s=3} & {625} & {111} & {21/5109/5637} \\
{256} & {N=11, s=2} & {600} & {103} & {22/5421/5528} \\
{512} & {CA0D2C1E [23]} & {625} & {275} & {31/3443/6602} \\
{512} & {[8]} & {162} & {33} & {188/29985/-} \\
{512} & {[8]} & {182} & {47} & {176/37138/-} \\
{512} & {[29]} & {--} & {--} & {1089/329868/4143} \\
{512} & {[18]} & {61.3} & {29} & {560/23511/11307} \\
{512} & {N=5, s=7} & {525} & {209} & {35/8044/10510} \\
{512} & {N=7, s=5} & {550} & {199} & {35/8124/10128} \\
{512} & {N=11, s=4} & {525} & {239} & {44/10498/12350} \\
{512} & {N=13, s=3} & {500} & {201} & {39/9266/11033} \\
{1024} & {CA0D2C1E [23]} & {625} & {545} & {61/6785/12108} \\
{1024} & {[11]} & {244} & {1372} & {9/5003/-} \\
{1024} & {N=5, s=13} & {525} & {401} & {65/14999/19272} \\
{1024} & {N=11, s=7} & {500} & {443} & {77/18282/21598} \\
{2048} & {CA0D2C1E [23]} & {625} & {1085} & {121/13487/22602} \\
{2048} & {[11]} & {246} & {2945} & {9/5964/-} \\
{2048} & {N=5, s=25} & {500} & {785} & {125/29182/35008} \\
{4096} & {CA0D2C1E [23]} & {625} & {2174} & {242/26978/44806} \\
{4096} & {[11]} & {175} & {2945} & {27/11898/-} \\
{4096} & {N=5, s=25} & {525} & {1553} & {245/58161/66740} \\
\end{longtable}
\endgroup
\begingroup
\small
\setlength{\tabcolsep}{3pt}
\setlength{\LTcapwidth}{\textwidth}
\begin{longtable}{lrrr}
\caption{Comparación de implementaciones recogida en el estudio de SQIsign: tiempo y producto área-tiempo (AT), en recurso por µs.}\label{tab:eval-sqisign-at}\\
\toprule
Bits & Diseño & Tiempo (µs) & AT: DSP/LUT/FF \\
\midrule
\endfirsthead
\toprule
Bits & Diseño & Tiempo (µs) & AT: DSP/LUT/FF \\
\midrule
\endhead
\midrule
\multicolumn{4}{r}{\small Continúa en la página siguiente.}\\
\endfoot
\bottomrule
\endlastfoot
{256} & {CA0D2C1E [23]} & {0.224} & {3.58/394/754} \\
{256} & {[8]} & {0.165} & {19.8/450/-} \\
{256} & {[8]} & {0.242} & {22.0/415/-} \\
{256} & {[29]} & {0.01755} & {5.072/1537/51} \\
{256} & {[5]} & {0.0852} & {21.13/805/-} \\
{256} & {[18]} & {0.334} & {40.1/2234/1724} \\
{256} & {[20]} & {0.0147} & {0/2762/-} \\
{256} & {N=3, s=6} & {0.178} & {3.20/738/914} \\
{256} & {N=5, s=4} & {0.181} & {3.62/872/1000} \\
{256} & {N=7, s=3} & {0.178} & {3.73/907/1000} \\
{256} & {N=11, s=2} & {0.172} & {3.78/931/949} \\
{512} & {CA0D2C1E [23]} & {0.440} & {13.6/1510/2900} \\
{512} & {[8]} & {0.204} & {38.4/6120/-} \\
{512} & {[8]} & {0.258} & {45.4/9580/-} \\
{512} & {[29]} & {0.02568} & {27.97/8471/106} \\
{512} & {[18]} & {0.472} & {264.3/11097/5336} \\
{512} & {N=5, s=7} & {0.398} & {13.9/3200/4180} \\
{512} & {N=7, s=5} & {0.362} & {12.7/2940/3660} \\
{512} & {N=11, s=4} & {0.455} & {20.0/4780/5620} \\
{512} & {N=13, s=3} & {0.365} & {14.3/3390/4030} \\
{1024} & {CA0D2C1E [23]} & {0.872} & {53.2/5920/10600} \\
{1024} & {[11]} & {5.63} & {50.7/28166/-} \\
{1024} & {N=5, s=13} & {0.764} & {49.6/11500/14700} \\
{1024} & {N=11, s=7} & {0.886} & {68.2/16200/19100} \\
{2048} & {CA0D2C1E [23]} & {1.74} & {210/23400/39000} \\
{2048} & {[11]} & {11.99} & {107.9/71508/-} \\
{2048} & {N=5, s=25} & {1.57} & {196/45800/54900} \\
{4096} & {CA0D2C1E [23]} & {3.48} & {842/93800/156000} \\
{4096} & {[11]} & {16.82} & {454.1/200124/-} \\
{4096} & {N=5, s=25} & {2.958} & {724/172000/197000} \\
\end{longtable}
\endgroup


\section{HAWK}

\subsection*{Ventajas}
\begin{itemize}
\item No utiliza operaciones en coma flotante lo que simplifica bastante
\item Implementación ya existente en FPGA basada en Dilithium (algo diferente)
\item De realizarse en HAWK, debería realizarse en SoC (PS+PL). Dificilmente viable puramente en HW.
\end{itemize}

\subsection*{Limitaciones}
\begin{itemize}
\item Sigue teniendo un algoritmo complejo con control complejo
\item No se ha encontrado ninguna implementación en HW ya existente
\item Implementación de NTT
\item Mayor complejidad que las multivariables
\end{itemize}

\section{Mirath}

\begingroup
\small
\setlength{\tabcolsep}{3pt}
\setlength{\LTcapwidth}{\textwidth}
\begin{longtable}{lrrrr}
\caption{Mirath: utilización de recursos para los niveles L1 y L5. Los bloques DSP son comunes a todos los niveles.}\label{tab:eval-mirath}\\
\toprule
Módulo & LUT L1/L5 & FF L1/L5 & DSP & BRAM L1/L5 \\
\midrule
\endfirsthead
\toprule
Módulo & LUT L1/L5 & FF L1/L5 & DSP & BRAM L1/L5 \\
\midrule
\endhead
\midrule
\multicolumn{5}{r}{\small Continúa en la página siguiente.}\\
\endfoot
\bottomrule
\endlastfoot
{Control principal (FSM)} & {958/1007} & {1003/1203} & {0} & {0} \\
{Control árbol GGM (FSM)} & {929/909} & {371/344} & {0} & {1/2} \\
{Motor SHA3} & {3757/3593} & {1978/1973} & {0} & {1} \\
{Motor AES/Rijndael} & {2073/4163} & {1278/2162} & {0} & {0} \\
{Núcleos aritméticos} & {7639/20699} & {4751/12477} & {0} & {34/73} \\
{Memoria de datos} & {0} & {0} & {0} & {4/16} \\
{Memoria de claves} & {0} & {0} & {0} & {2/4} \\
{FIFO de compromisos} & {44/44} & {64/64} & {0} & {0/0} \\
{Total} & {15400/30415} & {9445/18223} & {0} & {42/96} \\
\end{longtable}
\endgroup


\subsection*{Ventajas}
\begin{itemize}
\item No hay aritmética compleja
\item Repetición de la misma operación (muy paralelizable)
\end{itemize}

\subsection*{Limitaciones}
\begin{itemize}
\item Realiza mucho hash (consumo de memoria)
\item El resto de las funcionalidades también hacen un uso intensivo de memoria
\item Tiene que realizarse un buen control
\item Pocas implementaciones en HW
\end{itemize}

\section{MQOM}

\subsection*{Ventajas}
\begin{itemize}
\item Operaciones lógicas cortas y sencillas
\item Poco uso de LUT para operaciones aritméticas
\item Puede acelerarse con DSP
\end{itemize}

\subsection*{Limitaciones}
\begin{itemize}
\item Combina tanto MPC como Multivariable (no es fácil de implementar)
\item Control complejo
\item Elevado uso de LUT por hashing
\item Poca implementación previa en FPGA
\end{itemize}

\section{PERK}

\subsection*{Ventajas}
\begin{itemize}
\item Operaciones lógicas simples
\item Bajo uso de bloques DSP
\item Muchas rondas (paralelizable)
\end{itemize}

\subsection*{Limitaciones}
\begin{itemize}
\item Sigue habiendo hashing (muchas LUT)
\item Mucho uso de memoria
\item Control complejo
\item Poca implementación existente en FPGA
\end{itemize}

\section{RYDE}

\subsection*{Ventajas}
\begin{itemize}
\item Muchas rondas (paralelizable)
\item Operaciones aritméticas
\end{itemize}

\subsection*{Limitaciones}
\begin{itemize}
\item Control complejo
\item Uso intensivo de recursos (LUT y BRAM)
\item Mucho hashing
\item Poca implementación existente en FPGA
\end{itemize}

\section{SDitH}

\begingroup
\small
\setlength{\tabcolsep}{3pt}
\setlength{\LTcapwidth}{\textwidth}
\begin{longtable}{lrrrr}
\caption{SDitH: utilización de recursos FPGA.}\label{tab:eval-sdith}\\
\toprule
Parámetros & LUT & FF & DSP & BRAM \\
\midrule
\endfirsthead
\toprule
Parámetros & LUT & FF & DSP & BRAM \\
\midrule
\endhead
\midrule
\multicolumn{5}{r}{\small Continúa en la página siguiente.}\\
\endfoot
\bottomrule
\endlastfoot
{SDitH-L1-GF256} & {16592} & {13778} & {0} & {164.5} \\
{SDitH-L1-GF251} & {17463} & {16386} & {196} & {164.5} \\
{SDitH-L3-GF256} & {22906} & {18381} & {0} & {356.0} \\
{SDitH-L3-GF251} & {24381} & {27941} & {382} & {356.0} \\
{SDitH-L5-GF256} & {32423} & {24992} & {0} & {521.5} \\
{SDitH-L5-GF251} & {34136} & {31406} & {472} & {521.5} \\
\end{longtable}
\endgroup


\subsection*{Ventajas}
\begin{itemize}
\item Pocas operaciones lógicas (muy MPC)
\item Operaciones con matrices
\item No usa isogenías ni NTT
\item Algo más de implementación en FPGA a pesar de ser MPC
\end{itemize}

\subsection*{Limitaciones}
\begin{itemize}
\item Mucho hashing
\item Uso intensivo de memoria
\item Control complejo
\end{itemize}

\section{MAYO}

\begingroup
\small
\setlength{\tabcolsep}{3pt}
\setlength{\LTcapwidth}{\textwidth}
\begin{longtable}{lrrr}
\caption{MAYO y UOV: comparación de plataformas y recursos (LUT/FF/DSP/BRAM, salvo área ASIC en mm²).}\label{tab:eval-mayo-recursos}\\
\toprule
Esquema & Trabajo & Plataforma & Recursos / área \\
\midrule
\endfirsthead
\toprule
Esquema & Trabajo & Plataforma & Recursos / área \\
\midrule
\endhead
\midrule
\multicolumn{4}{r}{\small Continúa en la página siguiente.}\\
\endfoot
\bottomrule
\endlastfoot
{UOV-I} & {[7] a} & {Artix7 90.8 MHz} & {28k/24k/2/66} \\
{UOV-I} & {[7] a} & {Artix7 90.8 MHz} & {34k/27k/2/66} \\
{UOV-I} & {[7] a} & {AC-A72 1.8 GHz} & {--} \\
{UOV-I} & {[5]} & {AC-M4 24 MHz} & {--} \\
{UOV-I} & {[5] b} & {IXG 6338 2 GHz} & {--} \\
{MAYO-I} & {[10] b} & {AC-M7 480 MHz} & {--} \\
{MAYO-I} & {[23]} & {Artix7 100 MHz} & {21k/13k/11/129} \\
{MAYO-I} & {Our e} & {Artix7 75 MHz} & {92k/33k/2/45.5} \\
{MAYO-I} & {Our e} & {Kintex7 100 MHz} & {94k/32k/2/45.5} \\
{MAYO-I} & {Our e} & {A-28nm 1.5 GHz} & {1.02 mm²} \\
{MAYO-I} & {Our f} & {Artix7 75 MHz} & {106k/38k/2/45.5} \\
{MAYO-I} & {Our f} & {Kintex7 100 MHz} & {112k/38k/2/45.5} \\
{MAYO-I} & {Our f} & {A-28nm 1.5 GHz} & {1.24 mm²} \\
{UOV-III} & {[7] c} & {Artix7 96 MHz} & {38k/29k/2/184.5} \\
{UOV-III} & {[7] c} & {Artix7 94.1 MHz} & {43k/35k/2/184.5} \\
{UOV-III} & {[7] c} & {AC-A72 1.8 GHz} & {--} \\
{MAYO-III} & {[5]} & {IXG 6338 2 GHz} & {--} \\
{MAYO-III} & {Our e} & {Kintex7 100 MHz} & {147k/52k/2/96} \\
{MAYO-III} & {Our e} & {A-28nm 1.5 GHz} & {1.71 mm²} \\
{MAYO-III} & {Our f} & {Kintex7 100 MHz} & {169k/60k/2/96} \\
{MAYO-III} & {Our f} & {A-28nm 1.5 GHz} & {1.97 mm²} \\
{UOV-V} & {[7] d} & {Artix7 82.5 MHz} & {77k/38k/2/356} \\
{UOV-V} & {[7] d} & {Artix7 92.6 MHz} & {83k/41k/2/359} \\
{UOV-V} & {[7] d} & {AC-A72 1.8 GHz} & {--} \\
{MAYO-V} & {[5]} & {IXG 6338 2 GHz} & {--} \\
{MAYO-V} & {Our e} & {Kintex7 100 MHz} & {215k/73k/2/194.5} \\
{MAYO-V} & {Our e} & {A-28nm 1.5 GHz} & {3.09 mm²} \\
{MAYO-V} & {Our f} & {Kintex7 100 MHz} & {246k/83k/2/194.5} \\
{MAYO-V} & {Our f} & {A-28nm 1.5 GHz} & {3.45 mm²} \\
\end{longtable}
\endgroup
\begingroup
\small
\setlength{\tabcolsep}{3pt}
\setlength{\LTcapwidth}{\textwidth}
\begin{longtable}{lrrrrr}
\caption{MAYO y UOV: latencias expresadas como ciclos/ms; k y M denotan miles y millones de ciclos.}\label{tab:eval-mayo-tiempos}\\
\toprule
Esquema & Trabajo & Plataforma & Claves & Firma & Verificación \\
\midrule
\endfirsthead
\toprule
Esquema & Trabajo & Plataforma & Claves & Firma & Verificación \\
\midrule
\endhead
\midrule
\multicolumn{6}{r}{\small Continúa en la página siguiente.}\\
\endfoot
\bottomrule
\endlastfoot
{UOV-I} & {[7] a} & {Artix7 90.8 MHz} & {11.0M/121.94} & {843k/9.29} & {284k/3.13} \\
{UOV-I} & {[7] a} & {Artix7 90.8 MHz} & {11.0M/121.91} & {779k/8.63} & {115k/1.27} \\
{UOV-I} & {[7] a} & {AC-A72 1.8 GHz} & {28.3M/15.73} & {13.3M/7.40} & {2.20M/1.25} \\
{UOV-I} & {[5]} & {AC-M4 24 MHz} & {5.24M/21.84} & {9.18M/38.25} & {4.88M/20.33} \\
{UOV-I} & {[5] b} & {IXG 6338 2 GHz} & {110k/0.05} & {460k/0.23} & {175k/0.08} \\
{MAYO-I} & {[10] b} & {AC-M7 480 MHz} & {--} & {42.9M/89.43} & {6.70M/11.88} \\
{MAYO-I} & {[23]} & {Artix7 100 MHz} & {996k/9.96} & {2.867M/28.67} & {--} \\
{MAYO-I} & {Our e} & {Artix7 75 MHz} & {29330/0.40} & {95191/1.28} & {30150/0.40} \\
{MAYO-I} & {Our e} & {Kintex7 100 MHz} & {29330/0.30} & {95191/0.96} & {30150/0.30} \\
{MAYO-I} & {Our e} & {A-28nm 1.5 GHz} & {29330/0.02} & {95191/0.06} & {30150/0.02} \\
{MAYO-I} & {Our f} & {Artix7 75 MHz} & {5799/0.08} & {31128/0.43} & {4587/0.05} \\
{MAYO-I} & {Our f} & {Kintex7 100 MHz} & {5799/0.06} & {31128/0.32} & {4587/0.04} \\
{MAYO-I} & {Our f} & {A-28nm 1.5 GHz} & {5799/0.004} & {31128/0.02} & {4587/0.003} \\
{UOV-III} & {[7] c} & {Artix7 96 MHz} & {16.7M/174.24} & {1.46M/15.26} & {823k/8.57} \\
{UOV-III} & {[7] c} & {Artix7 94.1 MHz} & {16.4M/174.94} & {1.19M/12.75} & {195k/2.07} \\
{UOV-III} & {[7] c} & {AC-A72 1.8 GHz} & {56.8M/31.56} & {34.5M/19.18} & {8.31M/4.62} \\
{MAYO-III} & {[5]} & {IXG 6338 2 GHz} & {508k/0.25} & {1.66M/0.83} & {610k/0.30} \\
{MAYO-III} & {Our e} & {Kintex7 100 MHz} & {65660/0.65} & {214751/2.15} & {66959/0.67} \\
{MAYO-III} & {Our e} & {A-28nm 1.5 GHz} & {65660/0.04} & {214751/0.14} & {66959/0.04} \\
{MAYO-III} & {Our f} & {Kintex7 100 MHz} & {11730/0.11} & {65673/0.65} & {9035/0.10} \\
{MAYO-III} & {Our f} & {A-28nm 1.5 GHz} & {11730/0.008} & {65673/0.04} & {9035/0.006} \\
{UOV-V} & {[7] d} & {Artix7 82.5 MHz} & {39.0M/473.53} & {3.30M/40.09} & {1.9M/23.29} \\
{UOV-V} & {[7] d} & {Artix7 92.6 MHz} & {38.4M/414.73} & {2.64M/28.56} & {364k/3.93} \\
{UOV-V} & {[7] d} & {AC-A72 1.8 GHz} & {291M/161.91} & {86.7M/48.18} & {18.6M/10.33} \\
{MAYO-V} & {[5]} & {IXG 6338 2 GHz} & {1.21M/0.60} & {4.14M/2.07} & {1.18M/0.59} \\
{MAYO-V} & {Our e} & {Kintex7 100 MHz} & {118027/1.19} & {384488/3.89} & {119603/1.19} \\
{MAYO-V} & {Our e} & {A-28nm 1.5 GHz} & {118027/0.08} & {384488/0.26} & {119603/0.08} \\
{MAYO-V} & {Our f} & {Kintex7 100 MHz} & {19897/0.20} & {111736/1.12} & {14713/0.14} \\
{MAYO-V} & {Our f} & {A-28nm 1.5 GHz} & {19897/0.01} & {111736/0.07} & {14713/0.009} \\
\end{longtable}
\endgroup
Las marcas a--f, las referencias entre corchetes y la denominación \textit{Our} se conservan de la tabla del documento de partida; no identifican resultados del presente trabajo.



\subsection*{Ventajas}
\begin{itemize}
\item Solo se hace hash una vez (poco uso de LUT)
\item Operaciones aritméticas
\item Altamente paralelizable
\item Trabajos en FPGA ya existentes
\end{itemize}

\subsection*{Limitaciones}
\begin{itemize}
\item Resolución del sistema de ecuaciones
\item Uso intensivo de memoria
\item Clave pública grande (coeficientes del sistema)
\end{itemize}

\section{QR-UOV}

\subsection*{Ventajas}
\begin{itemize}
\item Operaciones aritméticas sencillas
\item Cálculo matricial
\item Poco hashing
\item Control asequible
\item Implementación de UOV ya existente
\item Reduce respecto a UOV el tamaño de la clave pública
\end{itemize}

\subsection*{Limitaciones}
\begin{itemize}
\item Mucho uso de memoria
\item Eliminación gaussiana (más reducida que UOV)
\end{itemize}

\section{SNOVA}

\subsection*{Ventajas}
\begin{itemize}
\item Operaciones aritméticas matriciales
\item Paralelizable
\item Repetitivo
\end{itemize}

\subsection*{Limitaciones}
\begin{itemize}
\item Mucho uso de memoria en matrices
\item Resolver el sistema es más complicado
\item Mayor control
\item Variables matriciales en lugar de escalares
\end{itemize}

\section{UOV}

\begingroup
\small
\setlength{\tabcolsep}{3pt}
\setlength{\LTcapwidth}{\textwidth}
\begin{longtable}{lrrrrr}
\caption{UOV: utilización de recursos.}\label{tab:eval-uov-recursos}\\
\toprule
Variante & Slices & LUT & FF & BRAM & DSP \\
\midrule
\endfirsthead
\toprule
Variante & Slices & LUT & FF & BRAM & DSP \\
\midrule
\endhead
\midrule
\multicolumn{6}{r}{\small Continúa en la página siguiente.}\\
\endfoot
\bottomrule
\endlastfoot
{uov-Ip-classic} & {12145} & {33221} & {24097} & {108.5} & {2} \\
{uov-Ip-pkc} & {12073} & {32134} & {22969} & {81} & {2} \\
{uov-Ip-pkc+skc} & {12106} & {32422} & {23262} & {48} & {2} \\
{uov-Is-classic} & {12860} & {44974} & {27433} & {140} & {2} \\
{uov-Is-pkc} & {11740} & {29385} & {25328} & {110} & {2} \\
{uov-Is-pkc+skc} & {11681} & {28947} & {24444} & {66} & {2} \\
{uov-III-pkc} & {17610} & {41761} & {31543} & {310.5} & {4} \\
{uov-III-pkc+skc} & {16574} & {38352} & {29446} & {184.5} & {4} \\
{uov-V-pkc+skc} & {27038} & {77352} & {38217} & {359} & {4} \\
\end{longtable}
\endgroup
\begingroup
\small
\setlength{\tabcolsep}{3pt}
\setlength{\LTcapwidth}{\textwidth}
\begin{longtable}{lrrrr}
\caption{UOV: ciclos de generación de claves, firma y verificación y frecuencia.}\label{tab:eval-uov-tiempos}\\
\toprule
Variante & Claves & Firma & Verificación & MHz \\
\midrule
\endfirsthead
\toprule
Variante & Claves & Firma & Verificación & MHz \\
\midrule
\endhead
\midrule
\multicolumn{5}{r}{\small Continúa en la página siguiente.}\\
\endfoot
\bottomrule
\endlastfoot
{uov-Ip-classic} & {3540971} & {7515} & {6435} & {93.5} \\
{uov-Ip-pkc} & {4170749} & {7515} & {192411} & {91.4} \\
{uov-Ip-pkc+skc} & {3807119} & {352621} & {192411} & {94.8} \\
{uov-Is-classic} & {9916182} & {13070} & {12986} & {92.2} \\
{uov-Is-pkc} & {11922375} & {13070} & {284379} & {94.8} \\
{uov-Is-pkc+skc} & {11027933} & {843885} & {284379} & {90.8} \\
{uov-III-pkc} & {18221241} & {19285} & {823108} & {97.5} \\
{uov-III-pkc+skc} & {16727607} & {1465182} & {823108} & {96.0} \\
{uov-V-pkc+skc} & {39066651} & {3308031} & {1921513} & {92.5} \\
\end{longtable}
\endgroup


\subsection*{Ventajas}
\begin{itemize}
\item Operaciones aritméticas
\item Solo se hace hash una vez
\item Previas implementaciones en FPGA
\end{itemize}

\subsection*{Limitaciones}
\begin{itemize}
\item Mucho uso de memoria
\item Clave pública grande
\item Requiere control
\end{itemize}

\section{FAEST}

\subsection*{Ventajas}
\begin{itemize}
\item Sin reducción gaussiana
\item Clave pública reducida
\item Uso de AES (ya ha sido implementado en FPGA)
\end{itemize}

\subsection*{Limitaciones}
\begin{itemize}
\item Alto consumo de LUT
\item Mayor complejidad que multivariante
\item Complejidad en el control
\item Poca computación
\end{itemize}
\section{Clasificación del estudio}
La tabla siguiente reproduce el orden de preferencia del documento de partida. La familia criptográfica se indica expresamente en lugar de utilizar la leyenda de colores original.

\begingroup
\small
\setlength{\tabcolsep}{3pt}
\setlength{\LTcapwidth}{\textwidth}
\begin{longtable}{lrr}
\caption{Clasificación de los candidatos en el estudio inicial.}\label{tab:eval-clasificacion}\\
\toprule
Posición & Candidato & Familia criptográfica \\
\midrule
\endfirsthead
\toprule
Posición & Candidato & Familia criptográfica \\
\midrule
\endhead
\midrule
\multicolumn{3}{r}{\small Continúa en la página siguiente.}\\
\endfoot
\bottomrule
\endlastfoot
{1º} & {QR-UOV} & {Multivariantes} \\
{2º} & {UOV} & {Multivariantes} \\
{3º} & {MAYO} & {Multivariantes} \\
{4º} & {SNOVA} & {Multivariantes} \\
{5º} & {CROSS} & {Basados en códigos} \\
{6º} & {HAWK} & {Retículos} \\
{7º} & {FAEST} & {Construcciones simétricas} \\
{8º} & {PERK} & {MPC-in-the-head} \\
{9º} & {SDitH} & {MPC-in-the-head} \\
{10º} & {RYDE} & {MPC-in-the-head} \\
{11º} & {Mirath} & {MPC-in-the-head} \\
{12º} & {LESS} & {Basados en códigos} \\
{13º} & {MQOM} & {MPC-in-the-head} \\
{14º} & {SQIsign} & {Isogenías} \\
\end{longtable}
\endgroup

